\chapter{Về Nỗ Lực}
\markboth{Về Nỗ Lực}{About Effort}

\section{Cố Gắng Mỗi Ngày}
\begin{truyen}{Cố Gắng Mỗi Ngày}{Try Every Day}
\chuhoa{N}{ỗ lực lớn không chỉ nằm ở những lúc bùng lên, mà ở việc em còn chịu cố một chút mỗi ngày.}
\textit{(Big effort is not only about dramatic moments, but about trying a little every day.)}

Điều đều đặn ấy nhìn có vẻ nhỏ, nhưng tích lại rất mạnh.
\textit{(That steady effort may look small, but it becomes powerful over time.)}

Người cố gắng mỗi ngày thường đi bền hơn người chỉ làm mạnh vài lần rồi tắt.
\textit{(Those who try every day often go farther than those who burn brightly only a few times.)}

Đó là một sự chăm chỉ rất thật.
\textit{(It is a very real kind of diligence.)}
\end{truyen}

\begin{baihoc}
Cố gắng mỗi ngày là cách làm cho một mục tiêu lớn trở nên có thể chạm tới hơn.
\textit{(Daily effort makes large goals more reachable.)}
\end{baihoc}

\begin{ghinhoanh}
\textbf{Short English to remember:}

Daily effort builds lasting progress.

\end{ghinhoanh}

\begin{tuhocnhanh}
1. Vì sao cố gắng mỗi ngày lại mạnh? \textit{Why is daily effort powerful?}

2. Nó khác gì với vài lần bùng lên? \textit{How is it different from occasional bursts?}

3. Em có thể cố điều gì mỗi ngày? \textit{What can you work on daily?}
\end{tuhocnhanh}
\ngancach

\section{Tiến Bộ Từng Chút}
\begin{truyen}{Tiến Bộ Từng Chút}{Improve Little by Little}
\chuhoa{T}{iến bộ thật nhiều khi đến rất nhỏ và rất yên.}
\textit{(Real progress often arrives very quietly and in small amounts.)}

Một chút chắc hơn, một chút đều hơn, một chút bớt sợ hơn.
\textit{(A little steadier, a little more regular, a little less afraid.)}

Nếu chỉ chờ những thay đổi lớn, học sinh sẽ dễ bỏ qua phần đẹp của việc lớn lên.
\textit{(If you wait only for big changes, you miss the beauty of gradual growth.)}

Tiến bộ từng chút vẫn là tiến bộ thật.
\textit{(Bit-by-bit progress is still real progress.)}
\end{truyen}

\begin{baihoc}
Đừng coi thường những bước nhỏ; chúng thường là cách thay đổi thật sự diễn ra.
\textit{(Do not underestimate small steps; they are often how real change happens.)}
\end{baihoc}

\begin{ghinhoanh}
\textbf{Short English to remember:}

Small growth is still real growth.

\end{ghinhoanh}

\begin{tuhocnhanh}
1. Vì sao tiến bộ từng chút vẫn quan trọng? \textit{Why does gradual progress still matter?}

2. Chờ thay đổi lớn dễ làm mình bỏ qua điều gì? \textit{What do we miss when we wait only for big change?}

3. Gần đây em đã tiến bộ từng chút ở đâu? \textit{Where have you improved little by little?}
\end{tuhocnhanh}
\ngancach

\section{Bền Bỉ Quan Trọng}
\begin{truyen}{Bền Bỉ Quan Trọng}{Persistence Matters}
\chuhoa{B}{ền bỉ là điều giữ nỗ lực không bị đứt giữa đường.}
\textit{(Persistence keeps effort from breaking halfway.)}

Có nhiều người có khả năng nhưng thiếu sự bền bỉ nên kết quả vẫn không đi xa.
\textit{(Many people have ability but lack persistence, so their results do not go far.)}

Ngược lại, người chịu bền lâu thường vượt lên chính mình nhiều hơn.
\textit{(By contrast, those who persist often outgrow themselves.)}

Bền bỉ vì thế không ồn ào, nhưng rất mạnh.
\textit{(Persistence is quiet, but powerful.)}
\end{truyen}

\begin{baihoc}
Trong nỗ lực, bền bỉ thường quan trọng hơn việc hăng hái ngắn hạn.
\textit{(In effort, persistence often matters more than short-lived enthusiasm.)}
\end{baihoc}

\begin{ghinhoanh}
\textbf{Short English to remember:}

Persistence keeps effort alive.

\end{ghinhoanh}

\begin{tuhocnhanh}
1. Vì sao bền bỉ quan trọng? \textit{Why does persistence matter?}

2. Thiếu bền bỉ dễ dẫn đến điều gì? \textit{What happens when persistence is missing?}

3. Em đang bền ở việc nào? \textit{Where are you being persistent?}
\end{tuhocnhanh}
\ngancach

\section{Không Có Đường Tắt}
\begin{truyen}{Không Có Đường Tắt}{There Is No Shortcut}
\chuhoa{N}{hiều học sinh mong có một cách thật nhanh để giỏi hơn mà không phải đi qua phần khó.}
\textit{(Many students hope for a quick way to improve without passing through difficulty.)}

Nhưng hầu hết năng lực thật đều cần thời gian, luyện tập và lặp lại.
\textit{(But most real abilities need time, practice, and repetition.)}

Đường tắt thường chỉ cho cảm giác nhanh trong chốc lát.
\textit{(Shortcuts usually feel quick only for a moment.)}

Về lâu dài, chúng hay để lại lỗ hổng lớn hơn.
\textit{(In the long run, they often leave bigger gaps.)}
\end{truyen}

\begin{baihoc}
Đi đúng đường chậm còn tốt hơn đi đường tắt rồi hụt nền.
\textit{(Taking the right road slowly is better than taking shortcuts and losing your foundation.)}
\end{baihoc}

\begin{ghinhoanh}
\textbf{Short English to remember:}

Real skill rarely comes by shortcut.

\end{ghinhoanh}

\begin{tuhocnhanh}
1. Vì sao không có đường tắt thật sự cho năng lực? \textit{Why is there no real shortcut to ability?}

2. Đường tắt thường để lại điều gì? \textit{What do shortcuts often leave behind?}

3. Em có đang tìm cách quá nhanh ở đâu không? \textit{Are you looking for a path that is too quick?}
\end{tuhocnhanh}
\ngancach

\section{Luyện Tập Nhiều}
\begin{truyen}{Luyện Tập Nhiều}{Practice Often}
\chuhoa{N}{ăng lực không lớn lên chỉ bằng việc hiểu một lần trong đầu.}
\textit{(Ability does not grow only by understanding something once in the mind.)}

Nó cần được dùng đi dùng lại để trở nên chắc tay và tự nhiên hơn.
\textit{(It needs to be used again and again to become steadier and more natural.)}

Luyện tập nhiều không phải chuyện hào nhoáng.
\textit{(Frequent practice is not glamorous.)}

Nhưng nó là phần âm thầm rất cần cho mọi tiến bộ thật.
\textit{(But it is the quiet part needed for real progress.)}
\end{truyen}

\begin{baihoc}
Điều em làm lặp lại nhiều lần sẽ dần trở thành sức của em.
\textit{(What you do repeatedly gradually becomes your strength.)}
\end{baihoc}

\begin{ghinhoanh}
\textbf{Short English to remember:}

Practice turns effort into skill.

\end{ghinhoanh}

\begin{tuhocnhanh}
1. Vì sao cần luyện tập nhiều? \textit{Why is frequent practice necessary?}

2. Điều gì xảy ra khi em lặp lại đủ lâu? \textit{What happens when you repeat long enough?}

3. Em cần luyện thêm ở kỹ năng nào? \textit{What skill do you need to practice more?}
\end{tuhocnhanh}
\ngancach

\section{Kiên Nhẫn Lâu Dài}
\begin{truyen}{Kiên Nhẫn Lâu Dài}{Long Patience}
\chuhoa{C}{ó những mục tiêu không thể nở ra trong vài ngày.}
\textit{(Some goals cannot bloom in just a few days.)}

Người biết kiên nhẫn lâu dài hiểu rằng nhiều điều tốt cần thời gian chín.
\textit{(Those with long patience understand that many good things need time to ripen.)}

Điều này giúp em bớt thất vọng mỗi khi kết quả chưa tới ngay.
\textit{(This helps you feel less discouraged when results do not come quickly.)}

Kiên nhẫn lâu dài là phần rất cần của nỗ lực bền.
\textit{(Long patience is an essential part of steady effort.)}
\end{truyen}

\begin{baihoc}
Muốn đi xa, em cần học cách chờ đợi mà vẫn tiếp tục làm phần của mình.
\textit{(If you want to go far, you must learn to wait while still doing your part.)}
\end{baihoc}

\begin{ghinhoanh}
\textbf{Short English to remember:}

Long patience supports long growth.

\end{ghinhoanh}

\begin{tuhocnhanh}
1. Vì sao cần kiên nhẫn lâu dài? \textit{Why is long patience necessary?}

2. Nó giúp em bớt điều gì? \textit{What does it help reduce?}

3. Em đang mong kết quả quá nhanh ở đâu? \textit{Where are you expecting results too quickly?}
\end{tuhocnhanh}
\ngancach

\section{Bắt Đầu Từ Việc Nhỏ}
\begin{truyen}{Bắt Đầu Từ Việc Nhỏ}{Start with Small Things}
\chuhoa{N}{hiều mục tiêu lớn làm học sinh sợ chỉ vì nhìn chúng quá to.}
\textit{(Many big goals frighten students simply because they look too large.)}

Trong khi đó, bắt đầu từ việc nhỏ giúp con đường có điểm vào rõ ràng hơn.
\textit{(Starting with small tasks gives the road a clearer entry point.)}

Một bước nhỏ nhưng thật thường đáng tin hơn một dự định lớn mà không bắt đầu.
\textit{(One real small step is often more trustworthy than a big plan that never begins.)}

Rất nhiều thành quả đẹp được ghép từ những việc nhỏ như thế.
\textit{(Many beautiful results are built from such small acts.)}
\end{truyen}

\begin{baihoc}
Đừng coi thường bước đầu nhỏ; nó thường là cách duy nhất để chuyện lớn bắt đầu.
\textit{(Do not underestimate the small first step; it is often the only way a big thing begins.)}
\end{baihoc}

\begin{ghinhoanh}
\textbf{Short English to remember:}

Small beginnings can carry big futures.

\end{ghinhoanh}

\begin{tuhocnhanh}
1. Vì sao nên bắt đầu từ việc nhỏ? \textit{Why should you start small?}

2. Điều gì đáng tin hơn một dự định lớn? \textit{What is more trustworthy than a big plan?}

3. Việc nhỏ nào em có thể bắt đầu hôm nay? \textit{What small thing can you start today?}
\end{tuhocnhanh}
\ngancach

\section{Làm Tốt Việc Hôm Nay}
\begin{truyen}{Làm Tốt Việc Hôm Nay}{Do Today's Work Well}
\chuhoa{T}{ương lai rất dễ làm học sinh lo, nhưng việc hôm nay mới là chỗ mình đang có thể làm được.}
\textit{(The future can make students anxious, but today's work is what we can actually do now.)}

Khi em làm tốt việc hôm nay, em đang âm thầm chuẩn bị cho ngày mai.
\textit{(When you do today's work well, you quietly prepare tomorrow.)}

Nhiều người mơ lớn nhưng lại lỏng ở việc trước mắt.
\textit{(Many people dream big but are careless with what is in front of them.)}

Nỗ lực thật luôn có chân đứng ở hiện tại.
\textit{(Real effort stands firmly in the present.)}
\end{truyen}

\begin{baihoc}
Làm tốt việc hôm nay là một trong những cách chắc nhất để đi tới ngày mai tốt hơn.
\textit{(Doing today's work well is one of the surest ways toward a better tomorrow.)}
\end{baihoc}

\begin{ghinhoanh}
\textbf{Short English to remember:}

Today's work shapes tomorrow.

\end{ghinhoanh}

\begin{tuhocnhanh}
1. Vì sao cần làm tốt việc hôm nay? \textit{Why should you do today's work well?}

2. Điều gì dễ khiến học sinh lỏng ở việc trước mắt? \textit{What makes students careless with present work?}

3. Việc hôm nay của em là gì? \textit{What is your work for today?}
\end{tuhocnhanh}
\ngancach

\section{Không Nản}
\begin{truyen}{Không Nản}{Do Not Lose Heart}
\chuhoa{N}{ản là điều rất thật, nhưng ở lại mãi trong nản thì việc học rất dễ dừng lại.}
\textit{(Discouragement is real, but staying there too long easily stops learning.)}

Người biết không nản không phải người không buồn.
\textit{(People who do not lose heart are not those who never feel low.)}

Họ chỉ là người vẫn còn chịu đi tiếp sau cảm giác ấy.
\textit{(They are people who continue after the feeling passes.)}

Đó là một sức mạnh rất thiết thực của nỗ lực.
\textit{(That is a very practical strength of effort.)}
\end{truyen}

\begin{baihoc}
Không nản không có nghĩa là không mệt; nó nghĩa là mệt mà vẫn chưa bỏ.
\textit{(Not losing heart does not mean never being tired; it means being tired without giving up.)}
\end{baihoc}

\begin{ghinhoanh}
\textbf{Short English to remember:}

Do not let discouragement make the final decision.

\end{ghinhoanh}

\begin{tuhocnhanh}
1. Vì sao không nản là một sức mạnh? \textit{Why is not losing heart a strength?}

2. Nó khác gì với không mệt? \textit{How is it different from never feeling tired?}

3. Em cần giữ mình đi tiếp ở chỗ nào? \textit{Where do you need to keep going?}
\end{tuhocnhanh}
\ngancach

\section{Đi Từng Bước}
\begin{truyen}{Đi Từng Bước}{Take It Step by Step}
\chuhoa{N}{hiều con đường dài trở nên khả thi hơn khi em không bắt mình nhảy quá xa cùng lúc.}
\textit{(Long roads become more manageable when you stop forcing yourself to leap too far at once.)}

Đi từng bước giúp mục tiêu lớn bớt đáng sợ.
\textit{(Step-by-step movement makes large goals less frightening.)}

Nó cũng giúp em thấy rõ mình đang tiến lên ở đâu.
\textit{(It also helps you see where you are progressing.)}

Trong nỗ lực, đi từng bước là một cách rất khôn ngoan.
\textit{(In effort, moving step by step is very wise.)}
\end{truyen}

\begin{baihoc}
Không cần nhảy quá xa để bắt đầu tiến bộ; nhiều khi chỉ cần một bước đủ thật.
\textit{(You do not need giant leaps to begin growing; often one honest step is enough.)}
\end{baihoc}

\begin{ghinhoanh}
\textbf{Short English to remember:}

Step-by-step effort is still real progress.

\end{ghinhoanh}

\begin{tuhocnhanh}
1. Vì sao đi từng bước lại khôn ngoan? \textit{Why is step-by-step effort wise?}

2. Nó giúp mục tiêu lớn thế nào? \textit{How does it help with large goals?}

3. Bước tiếp theo của em là gì? \textit{What is your next step?}
\end{tuhocnhanh}
\ngancach